简谐振子 \(m\ddot x+kx=0\) 的能量 \(E=\frac12 m\dot x^2+\frac12 kx^2\) 守恒。这不是偶然的——时间平移不变性（方程不含 \(t\)）通过 Noether 定理直接给出能量守恒。对称性不是比喻，是可计算的守恒律来源。

前面的 Part 分别处理了数学的各个分支：分析、代数、概率、优化、学习。这个单元换一个视角：不看分支，看结构。许多现象——从物理定律到数据分布——都呈现出低维的骨架。数学的价值不只是计算这些现象，而是揭示这些结构为什么存在、什么时候稳定、什么时候失效。

本单元回答四个问题：

\begin{enumerate}
\tightlist
\item
  对称性怎样产生守恒律：Noether 定理把连续对称性与守恒量接起来，从时间平移到能量守恒，从空间旋转到角动量守恒。
\item
  低维结构为什么普遍存在：流形假设、充分统计量、随机矩阵的普适性，都指向同一个现象——高维现象往往由少数自由度决定。
\item
  信息几何怎样统一统计与优化：统计流形上的 Fisher 度量、自然梯度、以及最大似然作为测地流，把估计问题变成几何问题。
\item
  因果结构的代数形式：潜在结果的代数、干预的算子、以及反事实的模态逻辑，把因果推断变成代数运算。
\end{enumerate}

\subsection{一个具体算例}

一维简谐振子：Lagrangian \(L=\frac12 m\dot x^2-\frac12 kx^2\)。时间平移 \(t\mapsto t+\epsilon\) 不改变 \(L\)（\(\partial L/\partial t=0\)），Noether 定理给出守恒量 \(E=\dot x\frac{\partial L}{\partial\dot x}-L=\frac12 m\dot x^2+\frac12 kx^2\)。

在统计中，指数族 \(\log p(x\mid\theta)=\theta^T T(x)-A(\theta)+\log h(x)\) 的充分统计量 \(T(X)\) 实现了数据的低维压缩：不管样本量多大，关于 \(\theta\) 的全部信息都在 \(\dim(\theta)\) 维的 \(T(X)\) 中。这是``结构决定可压缩性''的可算版本。

阅读本单元前，建议先熟悉 Part 04 的群作用、Part 11 的 Fisher 信息、Part 12 的流形与信息几何。这些前置概念会在本单元中作为``结构''的具体例子出现。

\subsection{怎么读这一章}

这四篇展示反复出现的结构：对称性限制变化，低维假设压缩自由度，信息几何修正参数尺度，因果代数区分观察与干预。阅读时同时寻找结构带来的力量与它依赖的假设。

\subsection{本章篇目}

以下篇目按概念依赖排列；若从具体问题进入，也可以先打开最接近当前困惑的一篇，再沿篇内链接回补。

\begin{itemize}
\tightlist
\item \hyperref[sec:12-Real-Analysis-Support/10-Essential-Structures/01]{``对称性怎样产生守恒律''}；
\item \hyperref[sec:12-Real-Analysis-Support/10-Essential-Structures/02]{``低维结构为什么普遍存在''}；
\item \hyperref[sec:12-Real-Analysis-Support/10-Essential-Structures/03]{``信息几何怎样统一统计与优化''}；
\item \hyperref[sec:12-Real-Analysis-Support/10-Essential-Structures/04]{``因果结构的代数形式''}。
\end{itemize}
